\documentclass[11pt]{article}
\usepackage[margin=25mm]{geometry}
\usepackage{fontspec}
\setmainfont{Latin Modern Roman}
\usepackage{amsmath,amssymb,mathtools}
\usepackage{unicode-math}
\setmathfont{Latin Modern Math}
\usepackage{xurl,hyperref}
\hypersetup{colorlinks=true,urlcolor=blue}
\setcounter{secnumdepth}{0}
\setlength{\parindent}{0pt}
\setlength{\parskip}{5pt}
\setlength{\emergencystretch}{3em}
\newcommand{\sref}[2]{\hyperlink{source:#1:#2}{[S#2]}}
\newcommand{\cataloguescope}[1]{\paragraph{Recorded target.}#1}
\newcommand{\R}{\mathbb R}
\newcommand{\C}{\mathbb C}
\newcommand{\N}{\mathbb N}
\newcommand{\Z}{\mathbb Z}
\newcommand{\Q}{\mathbb Q}
\newcommand{\F}{\mathbb F}
\newcommand{\eps}{\varepsilon}
\begin{document}
\section{Normal Hessian blocks, an exact four-body index jump, and stability obstructions}\label{statement}

\subsection{Definitions and mathematical statement}

\paragraph{Shared notation.}$\mathbb N=\{1,2,\ldots\}$; the standard number fields have their usual meanings. All additional conventions are stated locally.


\paragraph{Definitions and notation.}In the Newtonian $n$-body problem, positions $q_i\in\mathbb R^d$ and masses $m_i>0$ obey $m_i\ddot q_i=\nabla_{q_i}U$, where $U(q)=\sum_{i<j}m_im_j/|q_i-q_j|$ and the gravitational constant is normalized to one. Here $T=\tfrac12\sum_i m_i|\dot q_i|^2$, $E=T-U$, and a collision has $q_i=q_j$. The center of mass can be fixed at zero; moment of inertia is $I=\sum_i m_i|q_i|^2$. A central configuration has $\nabla_{q_i}U=-\lambda m_iq_i$ in center-of-mass coordinates for a common scalar $\lambda$. After fixing $I=1$, it is a critical point of $U$ on the collision-free inertia ellipsoid. Congruent configurations and scales are factored out for a finiteness count. A relative equilibrium is a rigidly rotating solution with time-independent shape. The Morse index of a constrained critical point is the number of negative Hessian directions, with multiplicity, after imposing the normalization and accounting for symmetry degeneracies. A planar configuration can be embedded in the spatial configuration space and acquire extra normal directions. Linear stability concerns the variational equations of the corresponding rotating motion.

\paragraph{Catalogue wording retained for traceability.}
Give a sharp upper bound for the Morse index of a nonplanar central configuration as a critical point of the positive Newtonian potential on normalized configuration space. Give a sharp lower bound for the index of a planar central configuration when viewed in the nonplanar configuration space. Determine the relation to stability of rigidly rotating periodic solutions; in particular, must a linearly stable relative equilibrium come from a minimum of that potential?
\paragraph{Scope and source qualification.}The source’s positive potential convention is made explicit. Planar versus ambient Hessians and Morse versus dynamical stability are distinct targets; the page-break artifact is removed without dropping any question. The catalogue title/status was paired with a machine-generated research attempt; that attempt is not peer-reviewed evidence of settlement. The statement covers the original source questions rather than only the title’s subcase.
\subsection{Short English statement}
Find sharp Morse-index bounds for central configurations and determine their relation to stability of rigid rotations.
\subsection{Research attempt}
\paragraph{Model, date and investigated scope.}
OpenAI GPT-6 Astra (configured model \texttt{gpt-6-astra}), ultra reasoning; generated 11 October 2026. Derived the exact out-of-plane constrained Hessian and its rotation zero modes, with a graph-eigenvalue index criterion. Sharp general index bounds and the stability/minimum relation remain unresolved.

\paragraph{Sources and literature check.}
Read the supplied complete definition, catalogue wording, scope qualifications and source list for this canonical ID. The original-source citations are inherited as supplied references; no fresh primary-source or novelty search was performed for this attempt. No current open-status claim is inferred from this calculation.

\paragraph{The normal Hessian block at a planar configuration.}Let \(q_i=(x_i,y_i,0)\) be a planar central configuration with multiplier \(\lambda>0\). An out-of-plane variation is \(\delta q_i=(0,0,z_i)\), satisfying center-of-mass constraint \(\sum m_iz_i=0\). It is automatically tangent to the inertia sphere. For each pair, differentiating \((r_{ij}^2+t^2(z_i-z_j)^2)^{-1/2}\) twice at zero gives \(-(z_i-z_j)^2/r_{ij}^3\). Thus the constrained Hessian is
\[Q_\perp(z)=\lambda\sum_im_iz_i^2-\sum_{i<j}\frac{m_im_j}{r_{ij}^3}(z_i-z_j)^2.
\]
Here the positive \(\lambda\) term comes from the Lagrange function \(U+\lambda I/2\), fixing the source's positive-potential convention.
\paragraph{Rotation zero modes and what can be counted.}Infinitesimal rotations tilting the plane give \(z_i=x_i\) and \(z_i=y_i\), both in the nullspace. This can also be checked directly: the central-configuration equations imply that the weighted graph Laplacian with edge weights \(m_im_j/r_{ij}^3\) sends each coordinate vector to \(\lambda\operatorname{diag}(m_i)\) times that vector. The normal Morse index is therefore the number of its generalized eigenvalues above \(\lambda\), after removing translation and the rotation zero modes. This is a computable block criterion once the configuration is known.
\paragraph{Gap.}No sharp uniform upper/lower index bound follows without information on those graph eigenvalues across all configurations. Dynamical stability also involves velocity/Coriolis terms in the rotating variational equation, not only the constrained potential Hessian. Consequently the stable-relative-equilibrium/minimum question is not answered by counting normal negative directions alone. Both sharp index questions and the stability implication remain unresolved; the first missing step is a global spectral bound for the relevant Hessian blocks and a separate reduced dynamical analysis.

\paragraph{Verification and adversarial audit.}
Checked pairwise second derivative, Lagrange multiplier sign, inertia tangent condition, center of mass and tilt zero modes. This was an LLM adversarial audit, not independent human verification; no theorem prover was used. Prompts A then B applied to this ID: Checked pairwise second derivative, Lagrange multiplier sign, inertia tangent condition, center of mass and tilt zero modes. Scope remains limited: Bound the relevant generalized graph eigenvalues uniformly across central configurations and analyze the full rotating linearization.

Coordinator checked the constrained normal Hessian sign, rotation null vectors and generalized graph-eigenvalue threshold; the dynamical gap remains explicit. This additional check was model-assisted, not independent human verification.

\paragraph{Tools, process and contributions.}
Python 3 and local shell used to read supplied files and serialize this manuscript; mathematics derived symbolically in the text, with no numerical experiment or formal proof assistant. Contribution prompt A was followed by prompt B. GPT-6 Astra with ultra reasoning drafted the research and performed its stated self-audit; any additional coordinator contribution is identified in the verification and audit. A GPT-6-based Codex coordinator assembled the manuscript and checked repository integration; its exact variant and reasoning configuration were not exposed. No human mathematical contribution, independent human proof check or formal proof certification is supplied. Complete derivations and the consulted source references are included here; private reasoning transcripts are not published.

\paragraph{Final status and first remaining gap.}
\textbf{UNRESOLVED}. Bound the relevant generalized graph eigenvalues uniformly across central configurations and analyze the full rotating linearization.

\section{Completion Estimate}
\noindent\textbf{COMPLETION: 0\%}
The calculation is an elementary benchmark or conditional reduction; it establishes no claimed new portion of the complete source target.

\subsection{Sources}
\begin{itemize}
\item[\hypertarget{source:statement:1}{S1}] ULAM.AI and the UnsolvedMath Contributors, UnsolvedMath, catalogue entry. Catalogue curation under CC BY 4.0; definitions and English exposition added locally.
\url{https://huggingface.co/datasets/ulamai/UnsolvedMath}.
\item[\hypertarget{source:statement:2}{S2}] American Institute of Mathematics, Variational Methods in Celestial Mechanics; notes by Kuo-Chang Chen, version 22 June 2003, source item 25.
\url{https://aimath.org/WWN/varcelest/varcelest.pdf}.
\end{itemize}
\label{end:statement}
\end{document}
